\begin{claim}{B.11d}{Proved}
Let $K$ be complete discretely valued of characteristic zero, with perfect
residue field of characteristic $p>0$, and let $U=\bar U\setminus D$ be
any hyperbolic curve. Good marked reduction is equivalent to the following
criterion in $p$-adic data of finite covers of $\Pi_U$:

There exist a finite \emph{unramified} extension $L/K$, a geometrically
connected finite Galois cover $V\to U_L$ of degree prime to $p$, and
$v\in V(L)$ such that its smooth proper compactification $C$ has genus
at least two, the finite map $C\to\bar U_L$ is ramified at every
geometric cusp, and the proper unipotent $p$-adic fundamental group of
$C$ based at $v$ is crystalline. Equivalently, every universal unipotent
representation $E^{\mathrm{et}}_{n,v}$ of $G_L$ for $C$ is crystalline.

In the forward direction the cover can have degree a power of any fixed
prime $\ell\ge5$, $\ell\ne p$. For $D=\varnothing$ the identity cover
is allowed. No semistability assumption on $U$ or $C$ is needed.
\end{claim}
\begin{proof}
We first identify the data attached to the original group. Such a cover
is given by a continuous quotient $q:E_{U_L}\twoheadrightarrow Q$
whose restriction to $\Pi_U$ is surjective. Put $H=\ker(q|_{\Pi_U})$.
The proper geometric group of its compactification is
\[
 \Pi_C=H\big/\overline{\langle H\cap\gamma I_c\gamma^{-1}:
                 c,\ \gamma\in\Pi_U\rangle}^{\mathrm{normal}},
                                                               \tag{B.11d.1}
\]
where $I_c$ ranges over the geometric cusp groups. Indeed a finite
etale cover of $V$ extends etale to $C$ exactly when its inertia at
each missing point vanishes, by normalization and the local parameter
criterion. Quotienting the arithmetic subgroup $\ker q$ by the same
normal subgroup gives $E_C$. The point $v$ defines its actual Galois
action; its $\Q_p$-unipotent completion is the data in the criterion.
We use these proper quotients of finite-index subgroups of $\Pi_U$,
with inherited arithmetic action and peripheral data, rather than
claiming direct crystallinity of $\Pi_U$ at every base point.

Suppose the pair has good reduction. A finite unramified extension
splits its etale marked divisor and supplies an integral point $b$ of
the smooth open model: a nonempty smooth special fibre has a point
over a finite separable residue extension, and henselian lifting
supplies $b$. We choose a finite $\ell$-group quotient in which every
cusp survives. If $\nu\ge2$, abelianization modulo $\ell$ suffices,
since the boundary classes have their sole relation that their sum is
zero and each class is nonzero. If $\nu=1$, then $g\ge1$, and the
geometric pro-$\ell$ group is free on $a_1,b_1,\ldots,a_g,b_g$,
with cusp $(\prod_i[a_i,b_i])^{-1}$. Map $a_1,b_1$ to the two
elementary generators of the Heisenberg group over $\F_\ell$, all
other generators to one; their commutator is a nonzero central element.
These presentations are SGA 1, XIII, Corollary 2.12 \citep{BC_SGA1}.

Intersecting all automorphic images of the kernel makes it
characteristic. Only finitely many occur: open subgroups of a fixed
index in a finitely generated profinite group are determined by
homomorphisms to a fixed finite symmetric group and a point stabilizer.
The new quotient $Q$ embeds in a finite product of the old finite
$\ell$-group, so remains an $\ell$-group and no cusp becomes trivial.
Every ramification index $e_c$ is at least $\ell$. The genus of the
compactification satisfies
\[
 2g(C)-2=|Q|\left(2g-2+\sum_{c\in D}(1-e_c^{-1})\right).
                                                               \tag{B.11d.2}
\]
For $g=0$, $\nu\ge3$ and $\ell\ge5$ bound the parentheses below by
$-2+3(1-1/5)>0$; for $g=1$, the nonempty ramification set makes them
positive; for $g\ge2$ positivity is automatic. Thus $g(C)\ge2$.
When $D$ is empty we instead use the identity cover, with $g\ge2$.

The integral section $b$ splits the arithmetic geometrically pro-$\ell$
extension. Smooth-pair comparison identifies its \emph{actual} action
with the action of its special-fibre point, so inertia is trivial.
The exact inputs are SGA 1, XIII, Proposition 4.3, Examples 4.4 and
4.5.3, with $\ell$ invertible on the trait \citep{BC_SGA1}.
The finite action on the characteristic quotient $Q$ therefore factors
through $G_k$. A further finite unramified extension kills it. Projection
to $Q$ is then a homomorphism of the semidirect product, trivial on
the section of $b$, and defines the desired constant-deck Galois cover
with an $L$-point $v$ above $b$.

The same quotient of the smooth open model's arithmetic group extends
the cover etale over that model. Normalize the proper model in its
function field. Ramification occurs only along the horizontal marks,
with indices powers of $\ell$. SGA 1, XIII, Propositions 5.1--5.2
give, after an etale change, the regular Kummer chart $z=w^{e_c}$.
Thus the normalization is smooth over the trait at the missing points
as well. The proper curve $C$ has good reduction, and $v$ is integral;
AIK supplies the crystalline representations
\citep[Theorems 1.6 and 1.9]{BC_AIK}.

Conversely suppose the criterion holds. B.11a gives good reduction
of the proper curve $C$, without semistability in advance. Its deck
automorphisms extend to its smooth proper model $\mathcal C/R_L$,
and their special action remains faithful by the finite unramified
stable-curve Isom scheme \citep[Theorem 1.11]{BC_DM}.
The projective quotient $\mathcal X=\mathcal C/Q$ exists by the
invariant graded ring of the canonical polarization. We prove locally
that it is smooth and that its reduced branch divisor is finite etale.

After strict henselization and completion, take a geometric special
point with stabilizer $H\subset Q$. Its tangent character is faithful.
Indeed a nontrivial finite-order automorphism of $\bar k[[z]]$ whose
order is prime to $p$ cannot have derivative one: its first nonzero
higher coefficient is multiplied by that order under iteration, and
cannot vanish. An automorphism identical on the completed local ring
is identical on the curve, contrary to faithfulness of the special
action. Thus $H$ embeds in $\bar k^*$ and is cyclic, of order $e$.

If $e>1$, a generator $h$ has special derivative $\zeta\ne1$.
The formal equation $h(z)-z=0$ has invertible derivative, so Hensel's
lemma gives a unique fixed section lifting the special point. Shift
$z$ to vanish on that section, lift $\zeta$ to the $e$th roots of
unity in the trait, and put
\[
              w=e^{-1}\sum_{i=0}^{e-1}\zeta^{-i}h^i(z).
\]
Its special linear coefficient is nonzero and $h(w)=\zeta w$.
The completed invariant ring is consequently
$R_L[[w]]^H=R_L[[w^e]]$. For $H=1$ the ring is simply $R_L[[w]]$.
To obtain the local ring of the global quotient, the completed
semilocal ring of a $Q$-orbit is a product of these local rings;
taking $Q$-invariants reduces it to $H$-invariants in one factor.
This proves smoothness of $\mathcal X$. At a ramified point the
branch locus is exactly the section $w^e=0$. No branch point appears
only in the special fibre, because every special fixed point lifted
uniquely, and no two branch points merge in this chart. Properness
then makes the full reduced branch divisor finite etale. Smoothness and this
divisor property descend through the strictly local base changes.

The generic quotient is $\bar U_L$, and its reduced branch divisor is
exactly $D_L$ by hypothesis. Thus $\mathcal X$ with that divisor is a
good marked model. Finally $L/K$ is unramified, so $I_L=I_K$ under
the usual common separable closure identification. Good marked
reduction over $L$ makes the pro-$\ell$ outer inertia trivial already
over $K$; Tamagawa's criterion gives good marked reduction over $R$
\citep[Theorem 5.3]{BC_Tamagawa}.
\end{proof}

This treats every marked type, including $(0,3)$, $(0,4)$ and $(1,1)$,
under the complete/perfect-residue hypotheses, through proper AIK on
a finite cover. A direct marked comparison on the original open
curve's unipotent group is not presumed. B.11f--B.12 give a separate
finite-cover $p$-divisible-group criterion over arbitrary residue fields.

\begin{claim}{B.11e}{Proved}
Let $R$ be henselian, of mixed characteristic, with perfect residue
field, and let $\widehat K$ be the fraction field of its completion.
Good marked reduction over $R$ is equivalent to the criterion B.11d
applied to $U_{\widehat K}$. Its Galois representations are attached
to the original arithmetic group via
\[
 G_{\widehat K}\simeq G_K,\qquad
 E_{U_{\widehat K}}\simeq E_U.
\]
In this formulation the auxiliary points in B.11d are points over
finite unramified extensions of $\widehat K$; no descent of those
points to $K$ is asserted.
\end{claim}
\begin{proof}
The natural continuous outer map $G_{\widehat K}\to G_K$ is an
isomorphism by the Krasner comparison, in exactly the henselian-DVR
scope of \citet[Lemma 3.1]{BC_MinTsuj2022}. Its map on residue Galois
groups is the identity, so it identifies inertia. Characteristic-zero
geometric base change identifies the geometric curve groups, including
the cusp subgroups, and the two arithmetic exact sequences therefore
identify $E_{U_{\widehat K}}$ with $E_U$.
Good marked reduction over $R$ implies it over $\widehat R$, hence
the finite-cover crystalline criterion by B.11d. Conversely that
criterion implies good marked reduction over $\widehat R$.
Thus the pro-$\ell$ outer inertia action is trivial there, for
$\ell\ne p$, and the identified inertia action is trivial over $K$.
Tamagawa's criterion over the original DVR gives good marked reduction
over $R$. This last argument avoids any unproved algebraization from
a formal model and retains the perfect residue field needed by AIK.
\end{proof}
